\subsection{Геометрия на плоскости}
\textbf{Треугольник} ($a,b,c$, $p=\frac{a+b+c}2$): $S=\sqrt{p(p{-}a)(p{-}b)(p{-}c)}=\frac12|\vec{AB}\times\vec{AC}|$; $R=\frac{abc}{4S}$, $r=\frac Sp$; медиана $m_a=\frac12\sqrt{2b^2{+}2c^2{-}a^2}$; биссектриса $\ell_a=\sqrt{bc\,(1-(\frac a{b+c})^2)}$ (делит $BC$ как $c{:}b$); теоремы синусов ($\frac a{\sin\alpha}=2R$) и косинусов. Инцентр $=\frac{aA+bB+cC}{a+b+c}$. $O$ -- центр описанной окружности: $H=A{+}B{+}C-2O$, центроид $G=\frac{2O+H}3$ (прямая Эйлера).
\par \textbf{Четырёхугольник} (стороны $a,b,c,d$, диагонали $e,f$, угол между диагоналями $\theta$, $F=b^2{+}d^2{-}a^2{-}c^2$): $4S=2ef\sin\theta=F\tan\theta=\sqrt{4e^2f^2-F^2}$. Вписанный: $ef=ac+bd$ (Птолемей), $S=\sqrt{(p{-}a)(p{-}b)(p{-}c)(p{-}d)}$.
\par \textbf{Многоугольник:} $2S=\sum p_i\times p_{i+1}$ ($>0$ против часовой); центр масс $\frac1{6S_{\pm}}\sum(p_i{+}p_{i+1})(p_i\times p_{i+1})$ при $S_\pm=\frac12\sum p_i\times p_{i+1}$.
\par \textbf{Окружности:} расстояние между центрами $d$; хорда пересечения на расстоянии $a=\frac{r_1^2-r_2^2+d^2}{2d}$ от центра 1, полудлина $h=\sqrt{r_1^2-a^2}$; радикальная ось: $2(c_2{-}c_1)\cdot X=|c_2|^2-|c_1|^2-r_2^2+r_1^2$. Касательные из $p$: пересечь окружность с окружностью $(p,\sqrt{|pc|^2-r^2})$.
\par \textbf{Углы:} угол между векторами лучше брать как $\mathrm{atan2}(u\times v,\ u\cdot v)$ (не $\arccos$). Поворот на $\varphi$: $(x\cos\varphi-y\sin\varphi,\ x\sin\varphi+y\cos\varphi)$. Сферические: $x=r\sin\theta\cos\phi,\ y=r\sin\theta\sin\phi,\ z=r\cos\theta$.
\par \textbf{3D:} объём тетраэдра $\frac16|(b{-}a)\cdot((c{-}a)\times(d{-}a))|$. Эйлер для выпуклого многогранника и связного плоского графа: $V-E+F=2$; планарный граф: $E\le3V-6$ (двудольный: $E\le2V-4$).
\par \textbf{Среднее и центр.} Среднее арифметическое минимизирует $\sum|x-p_i|^2$: $f(q)=f(\bar p)+n|q-\bar p|^2$ (для суммы расстояний -- медиана). Триангуляция выпуклого $n$-угольника: двойственный граф -- дерево из $n{-}2$ вершин со степенями $\le3$; триангуляций $C_{n-2}$.
\par \textbf{Делоне:} вне описанной окружности каждого треугольника нет других точек; единственна, если нет 4 точек на одной окружности; максимизирует минимальный угол. Вороной -- пересечение полуплоскостей серединных перпендикуляров к соседям по Делоне.
\subsection{Теоремы о графах}
\textbf{Паросочетания.} Кёниг (двудольный): макс. паросочетание $=$ мин. вершинное покрытие. Покрытие по Кёнигу: DFS из свободных вершин левой доли по неприкрытым и по паросочетанию обратно; $L^-\cup R^+$ -- покрытие, $L^+\cup R^-$ -- независимое множество. Галлаи: $\alpha+\beta=n$, $\nu+\rho=n$ (без изолированных). Холл: паросочетание, насыщающее $L$, $\iff|N(S)|\ge|S|\ \forall S\subseteq L$. Татт--Берж: $\nu(G)=\min_{U}\frac12(|V|+|U|-\mathrm{odd}(G-U))$.
\par \textbf{Дилворт:} макс. антицепь $=$ мин. покрытие цепями $=n-$ макс. паросочетание в двудольном графе строгого порядка (нужно \emph{транзитивное замыкание}). Антицепь: берём вершины, ни одна из двух копий которых не в покрытии Кёнига. Мин. покрытие DAG непересекающимися путями $=n-$ паросочетание (без замыкания).
\par \textbf{Степенные последовательности.} Эрдёш--Галлаи: $d_1\ge\dots\ge d_n$ -- степени простого графа $\iff\sum d_i$ чётна и $\forall k:\ \sum_{i\le k}d_i\le k(k{-}1)+\sum_{i>k}\min(d_i,k)$. Хавел--Хакими: убрать наибольшее $d_1$, вычесть 1 у следующих $d_1$ больших.
\par \textbf{Раскраска.} Хроматический многочлен: путь/дерево $t(t{-}1)^{n-1}$, цикл $(t{-}1)^n+(-1)^n(t{-}1)$, клика $t(t{-}1)\cdots(t{-}n{+}1)$; $P(G)=P(G{-}e)-P(G/e)$. Брукс: $\chi\le\Delta$ (кроме клик и нечётных циклов), Визинг: рёберная $\le\Delta{+}1$ (двудольный -- $\Delta$). Мантель: без треугольников $\le\lfloor n^2/4\rfloor$ рёбер.
\par \textbf{Потоки.} Макс-флоу $=$ мин-разрез. Поток с нижними границами: ребро $u{\to}v$ с $[l,c]$ заменяем ребром ёмкости $c{-}l$ и сдвигаем баланс ($\mathrm{ex}[v]{+}{=}l,\ \mathrm{ex}[u]{-}{=}l$); добавляем $S{\to}v$ ёмкости $\mathrm{ex}[v]$ при $\mathrm{ex}[v]>0$ и $v{\to}T$ при $<0$; допустимая циркуляция есть $\iff$ макс-флоу насыщает все рёбра из $S$. Менгер: число рёберно-непересекающихся $s$--$t$ путей $=$ мин. разрез. Гомори--Ху: $n{-}1$ потоков строят дерево всех мин. разрезов. Стабильные паросочетания: заявки в порядке предпочтений, отвергнутые идут дальше.
\par \textbf{Матроиды.} Пересечение двух матроидов: граф замен $A\setminus u\cup v\in I_1$ (рёбра $u{\to}v$), $\in I_2$ ($v{\to}u$); источники $\{z: A{\cup}z\in I_1\}$, стоки $\{z:A{\cup}z\in I_2\}$; кратчайший путь $S{\to}T$ увеличивает $A$ (взвешенный -- кратчайший по весу, среди них по числу рёбер).
\par \textbf{Цепи Маркова на графе.} Случайное блуждание по связному неориентированному неразобщённому графу: стационарное $\pi_v\propto\deg v$. Поглощающая цепь: $N=(I-Q)^{-1}$, ожидаемое время $=N\mathbf1$, вероятности поглощения $=NR$.
